\section{Limitations}

The study has four limitations. First, evaluation uses six categories, as
required by the course handout, rather than all twelve VisA categories.
Second, AnomalyDINO processes one image at a time in the official code. This
preserves exact reproduction but leaves substantial GPU capacity unused during
image decoding, PCA masking, and CPU nearest-neighbor search. Third,
patch-level masks are evaluated on the native $32\times32$ token grid rather
than after full-resolution TIFF upsampling, reflecting both the model's true
localization granularity and the server's tight disk budget. Fourth, the
multi-layer choice was fixed before test evaluation and is therefore not
selected with test labels, but it is still a single architectural hypothesis;
other layer pairs or learned scale weights may perform better.

The template baseline is intentionally simple. More powerful alignment,
foreground segmentation, or SSIM could improve it, but would also make the
comparison less interpretable. Its purpose is to expose the gap between raw
pixel correspondence and foundation-model patch retrieval, not to maximize
performance.
